\begin{ejercicio}
    Sean $M$ y $N$ módulos semisimples con descomposiciones como suma directa de submódulos simples $M = S_1\oplus \ldots\oplus S_t$ y $N=T_1\oplus \ldots\oplus T_s$. Supongamos que $S_i$ no es isomorfo a $T_j$ para todo $i = 1,\ldots,t,j=1,\ldots,s$. Demostrar que todo homomorfismo de módulos de $M$ a $N$ es cero.
\end{ejercicio}

\begin{ejercicio}\label{ej:39}
    Sea $M$ un módulo semisimple de dimensión finita con estructura $(\Sigma_1,n_1), \ldots, (\Sigma_t,n_t)$. Si $N$ es un submódulo de $M$, demostrar que su estructura es $(\Sigma_1,m_1), \ldots, (\Sigma_t,m_t)$ para ciertos $m_j\leq n_j$ (admitimos que $m_j=0$ significa que $\Sigma_j$ no aparece en la estructura de $N$).
\end{ejercicio}

\begin{ejercicio}
    Establecer un enunciado análogo al del Ejercicio~\ref{ej:39} para cada cociente de $M$.
\end{ejercicio}

\begin{ejercicio}
    Dada una $K-$álgebra $A$, demostrar que la aplicación $\rho:A\to {(End_A(A))}^{op}$ definida por $\rho(a)(a') = a'a$ es un isomorfismo de $K-$álgebras.
\end{ejercicio}

\begin{ejerciciostar}
    Sea $B$ un álgebra. Demostrar que la aplicación que asigna a cada matriz su traspuesta da un isomorfismo de álgebras ${(M_n(B))}^{op} \ong M_n(B^{op})$.
\end{ejerciciostar}

\begin{ejercicio}
    Sea $\varphi:R\to S$ un isomorfismo de $K-$álgebras e $I,J$ ideales a izquierda de $R$. Demsotrar que $\varphi(I),\varphi(J)$ son ideales a izquierda de $S$ y que dado cualquier homomorfismo de $R-$módulos $f:I\ŧo J$, la aplicación $\hat{f}:\varphi(I)\to \varphi(J)$ definida por $\hat{f}(y) = \varphi f \varphi^{-1}(y)$ para $y\in \varphi(I)$ es un homomorfismo de $S-$módulos.
\end{ejercicio}

\begin{ejercicio}
    Sean $R_1,\ldots,R_n$ $K-$álgebras y $R=R_1\times \ldots\times R_n$. Los ideales por la izquierda de $R$ son de la forma $I_1\times \ldots\times I_n$, con $I_i$ ideal a izquierda de $R_i$ para $i = 1,\ldots,n$. Análoga descripción tienen los ideales de $R$.
\end{ejercicio}

\begin{ejerciciostar}
    Sea $A$ un álgebra simple finito-dimensional. Demostrar que $R=M_n(A)$ es un álgebra simple de dimensión finita. Demostrar asimismo que si $\Sigma$ es un $A-$módulo simple y $M$ es un $R-$módulo simple, entonces $End_A(\Sigma)$ y $End_R(M)$ son álgebras de división isomorfas.
\end{ejerciciostar}

\begin{ejercicio}
    Sea $R$ un anillo y $\{e_1,\ldots,e_m\}$ un CCIO centrales de $R$. Sea $S$ un $R-$módulo simple. Demostrar que existe un único $j\in \{1,\ldots,m\}$ tal que $e_jS\neq \{0\}$. Deducir que, para $s\in S$; se tiene que $e_js = s$ y $e_is = 0$ para $i\neq j$.
\end{ejercicio}
